\section{An \texorpdfstring{$\ell^2$}{l2} short-shift transfer for the full dual ball}
\label{sec:l2-transfer}

We now change the shift quantifier.  The underlying prime theorem is the
almost-all short-interval result of \citet{GuthMaynard2026}; the point here
is that its exceptional set yields a small vector in the full dual
coefficient space.

For integer $h_0$ and parameters satisfying
\eqref{eq:short-shift-range} and
\eqref{eq:translation-geometric-range}, define, on the product-center
interval in \eqref{eq:W-support},
\begin{equation}
 \Delta_{X,H,h_0}(r)=
 \frac1H\sum_{q\in\Z}G(q/H)\eta(r+h_0+q).
 \label{eq:translated-residual-vector}
\end{equation}

\begin{theorem}[Translated prime residual in $\ell^2$]
\label{thm:translated-residual-l2}
Fix $0<\eps<13/15$.  For every $A>0$, uniformly in the stated geometric
range,
\begin{equation}
 \left(
 \sum_{\alpha X\leq r\leq\beta X}
 |\Delta_{X,H,h_0}(r)|^2
 \right)^{1/2}
 \ll_{A,\eps,W,G}X^{1/2}(\log X)^{-A}.
 \label{eq:translated-residual-l2}
\end{equation}
\end{theorem}

\begin{proof}
We give the smoothing argument because $G$ may have two-sided support.
First suppose $H\leq X^{0.98}$ and put
\[
 L=\left\lceil\exp\{(\log X)^{1/8}\}\right\rceil.
\]
Partition a fixed interval containing $\supp G$ into $L$ equal pieces.
The corresponding integer shift blocks have length
$H/L+O(1)\geq X^{2/15+\eps/2}$ for large $X$, and at most
$O(X^{0.98})$.  Apply Corollary~1.4 of \citet{GuthMaynard2026} to each
forward block.  Negative shifts merely change the starting point, which
remains comparable with $X$ by
\eqref{eq:translation-geometric-range}.  Taking the union over $O(L)$
blocks gives a set
$\mathfrak E\subset[\alpha X,\beta X]\cap\Z$ with
\begin{equation}
 |\mathfrak E|\ll_{\eps,W,G}
 X\exp\{-c_{\eps,G}(\log X)^{1/4}\}.
 \label{eq:l2-exception-size}
\end{equation}
Outside this set, the prime number theorem holds on every block with a
stretched-exponential relative error.  Freeze $G$ on each block.  Its
$C^1$ variation contributes $O_G(H\log X/L)$ by the elementary bound
$\Lambda\leq\log(3X)$, which is
$O_{B,G}(H(\log X)^{-B})$ for every fixed $B$.  Partial summation and the
negligible proper prime powers therefore give, outside $\mathfrak E$,
\begin{equation}
 |\Delta_{X,H,h_0}(r)|\ll_{B,\eps,G}(\log X)^{-B}
 \label{eq:l2-regular-pointwise}
\end{equation}
for every $B>0$.

If $H>X^{0.98}$, partition the shift range into consecutive blocks of
length between $X^{0.97}$ and $2X^{0.97}$.  Corollary~1.3 of
\citet{GuthMaynard2026}, whose all-interval range begins below this
length, applies to every block.  Freezing $G$ now costs
$O_G(X^{0.97}\log X)=O_{B,G}(H(\log X)^{-B})$.  Thus
\eqref{eq:l2-regular-pointwise} holds at every center in the long-window
range, and we take $\mathfrak E=\varnothing$.  This also explains why the
claimed range extends beyond the $X^{0.99}$ upper endpoint of the
almost-all corollary.

At an exceptional center, the elementary upper-bound sieve, or simply the
pointwise bound for $\Lambda$, gives
\begin{equation}
 |\Delta_{X,H,h_0}(r)|\ll_G\log(3X).
 \label{eq:l2-exception-pointwise}
\end{equation}
Choose $B>A+2$.  Squaring \eqref{eq:l2-regular-pointwise} and summing over
$O(X)$ regular centers gives $O(X(\log X)^{-2A})$.  Equations
\eqref{eq:l2-exception-size}--\eqref{eq:l2-exception-pointwise} give a
smaller bound on the exceptional set, because the stretched exponential
absorbs every fixed logarithmic power.  Taking square roots proves the
theorem.  The $O_G(H^{-1})$ Riemann-sum error in the constant part of
$\eta=\Lambda-1$ is absorbed since $H$ is a fixed positive power of $X$.
\end{proof}

\begin{corollary}[Abstract coefficient-space transfer]
\label{cor:abstract-coefficient-transfer}
Suppose $c_X(r)$ is supported on $[\alpha X,\beta X]$ and, for some fixed
$B\geq0$,
\begin{equation}
\sum_r|c_X(r)|^2\ll X(\log X)^B.
\label{eq:coefficient-l2-hypothesis}
\end{equation}
Assume that the implied constant is independent of $X,H,h_0$.
Then, for every $A>0$,
\begin{equation}
 \frac1H\sum_qG(q/H)
 \sum_rc_X(r)\eta(r+h_0+q)
 \ll_{A,B,\eps,W,G,c}X(\log X)^{-A}.
 \label{eq:abstract-coefficient-transfer}
\end{equation}
\end{corollary}

\begin{proof}
Cauchy--Schwarz and \cref{thm:translated-residual-l2}, requested with
$A+B/2+2$ in place of $A$, give the result.
\end{proof}

We now verify that every bounded full-coordinate mask belongs to the same
coefficient ball.  Group \eqref{eq:masked-residual} by $r=abk^2$:
\begin{equation}
 \cZ_{h,D}^{\omega}(X)
 =\sum_rC_{\omega,D}(r)\eta(r+h),
 \label{eq:masked-product-grouping}
\end{equation}
where
\begin{equation}
 C_{\omega,D}(r)=W(r/X)
 \sum_{\substack{abk^2=r,\ (a,b)=1\\ak>D}}
 \omega_{X,D}(a,b,k)u(ak).
 \label{eq:masked-product-coefficient}
\end{equation}

\begin{lemma}[Uniform bounded-mask ledger]
\label{lem:bounded-mask-ledger}
Uniformly in $D\geq1$ and every mask satisfying
\eqref{eq:bounded-mask},
\begin{align}
 |C_{\omega,D}(r)|
 &\leq\|W\|_\infty\bigl(\log r+\tau(r)\bigr),
 \label{eq:mask-pointwise-ledger}\\
 \sum_r|C_{\omega,D}(r)|^2
 &\ll_WX(\log X)^3.
 \label{eq:mask-l2-ledger}
\end{align}
\end{lemma}

\begin{proof}
The bijection \eqref{eq:gcd-coordinates} sends the triples with
$abk^2=r$ to the factor pairs $(n,m)$ with $nm=r$, and $ak=n$.
Therefore
\[
 |C_{\omega,D}(r)|
 \leq\|W\|_\infty\sum_{n\mid r}|\Lambda(n)-1|
 \leq\|W\|_\infty
 \left(\sum_{n\mid r}\Lambda(n)+\tau(r)\right).
\]
The divisor identity $\sum_{n\mid r}\Lambda(n)=\log r$ proves
\eqref{eq:mask-pointwise-ledger}.  The classical second moment
\citep{MontgomeryVaughan2007}
\[
 \sum_{r\leq y}\tau(r)^2\ll y(\log(2y))^3
\]
then proves \eqref{eq:mask-l2-ledger}.
\end{proof}

\begin{theorem}[Uniform full-label mask transfer]
\label{thm:uniform-mask-transfer}
Fix $0<\eps<13/15$.  For every $A>0$, uniformly in $D\geq1$, every mask
\eqref{eq:bounded-mask}, and the range of
\cref{thm:translated-residual-l2},
\begin{equation}
 \frac1H\sum_qG(q/H)
 \cZ_{h_0+q,D}^{\omega}(X)
 \ll_{A,\eps,W,G}X(\log X)^{-A}.
 \label{eq:uniform-mask-transfer}
\end{equation}
The implied constant is independent of the mask.
\end{theorem}

\begin{proof}
Apply \cref{cor:abstract-coefficient-transfer} to
\eqref{eq:masked-product-coefficient}, using
\eqref{eq:mask-l2-ledger}.
\end{proof}

There is an equivalent raw-target formulation.  Put
\begin{align}
 \cT_{h,D}^{\omega}(X)
 &=\sum_{\substack{a,b,k\geq1,\ (a,b)=1\\ak>D}}
 \omega_{X,D}(a,b,k)u(ak)
 \Lambda(abk^2+h)W(abk^2/X),
 \label{eq:raw-masked-transform}\\
 \cM_D^{\omega}(X)
 &=\sum_{\substack{a,b,k\geq1,\ (a,b)=1\\ak>D}}
 \omega_{X,D}(a,b,k)u(ak)W(abk^2/X).
 \label{eq:raw-mask-model}
\end{align}

\begin{corollary}[Uniform raw mask model]
\label{cor:uniform-raw-mask-model}
Under the hypotheses of \cref{thm:uniform-mask-transfer},
\begin{equation}
 \frac1H\sum_qG(q/H)\cT_{h_0+q,D}^{\omega}(X)
 =I_0(G)\cM_D^{\omega}(X)
 +O_{A,\eps,W,G}\bigl(X(\log X)^{-A}\bigr).
 \label{eq:uniform-raw-mask-model}
\end{equation}
\end{corollary}

\begin{proof}
Write $\Lambda=1+\eta$ and use
\cref{thm:uniform-mask-transfer}.  Euler summation gives
$H^{-1}\sum_qG(q/H)=I_0(G)+O_G(H^{-1})$.  The pointwise ledger and the
mean order of $\tau$ give
$|\cM_D^\omega|\ll_WX\log(2X)$, so the Riemann-sum error is absorbed in the
displayed remainder.
\end{proof}

\begin{remark}[Uniform masks are not uniform target adaptation]
The theorem permits a different mask for every $X,D,H$, including masks
chosen after inspecting the source labels.  It does not permit
$\omega=\omega_h$ to inspect each translated target.  Such a family could
retain precisely the prime targets and would make the assertion circular.
\end{remark}
